\documentclass[10pt,a4paper]{amsart}
\usepackage[margin=19mm]{geometry}
\usepackage{amsmath,amssymb,mathtools}
\usepackage[scr=rsfs,cal=euler]{mathalfa}
\usepackage{xfrac}
\usepackage[unicode,hidelinks,bookmarks=false]{hyperref}
\newcommand{\ie}{i.e.\ }
\newcommand{\cf}{cf.\ }
\newcommand{\resp}{respectively\ }
\newcommand{\Subsection}{\subsection}
\providecommand{\nobreakdash}{\nobreak\hbox{-}}
\setlength{\emergencystretch}{2em}
\multlinegap0pt
\usepackage{bm}
\usepackage{bbm}
\usepackage{dsfont}
\usepackage{xspace}
\usepackage{cancel}
\usepackage{mathtools}
\usepackage{amsthm}
\makeatletter
\@ifundefined{theorem}{\newtheorem{theorem}{Theorem}}{}
\makeatother
\makeatletter
\expandafter\ifx\csname acknowledgement\endcsname\relax
\newenvironment{acknowledgement}[1][Acknowledgement
]{\begin{trivlist} \item[\hskip \labelsep {\bfseries
#1}]}{\end{trivlist}}
\fi
\expandafter\ifx\csname example\endcsname\relax
\newenvironment{example}[1][Example]{\begin{trivlist}
\item[\hskip \labelsep {\bfseries #1}]}{\end{trivlist}}
\fi
\makeatother
\title[Uniqueness of Inverse Spectral Problem of Non-Local Sturm-Li]{Uniqueness of Inverse Spectral Problem of Non-Local Sturm-Liouville Operators on Star Graph}
\author{Lung-Hui Chen}
\date{Source: arXiv:2508.11102v1 (CC BY 4.0)}
\begin{document}
\maketitle

\noindent\emph{Source and licence.} Uniqueness of Inverse Spectral Problem of Non-Local Sturm-Liouville Operators on Star Graph. Version arXiv:2508.11102v1 (CC BY 4.0), distributed under the Creative Commons Attribution 4.0 International licence, as recorded in the arXiv licence metadata for this exact version and shown on the abstract page https://arxiv.org/abs/2508.11102v1. Full rights notice: licenses/arxiv-2508.11102-CC-BY-4.0.txt.\par\smallskip

\noindent\textit{Editorial scope.} Counted unit: the Theorem whose source label is \texttt{None} in the article (source0.tex lines 282--286, source label \texttt{None}), delivered together with the article's own complete original proof of it (source0.tex lines 287--302, one \texttt{proof} environment of 16 lines). No mathematics is written, repaired or re-derived: statement and proof are the article's own text line for line. Required context retained uncounted: 1.1 (lines 62--64); 1.2 (lines 67--69); 1.3 (lines 71--73); Phi (lines 111--114); T31 (lines 165--167); 3.4 (lines 211--215); 3.6 (lines 219--223); 3.7 (lines 226--230); 3.11 (lines 232--235). Every definition, display and label of those lines is retained unchanged. Omitted from this package and not redistributed: 1--61, 65--66, 70--70, 74--110, 115--164, 168--210, 216--218, 224--225, 231--231, 236--281, 303--362. Nothing inside a retained excerpt is omitted or altered: every retained body is delivered exactly as the article's lines are written, so no omission receipt of any kind accompanies this package. Citations: the retained excerpts carry no citation command at all, so no bibliography entry of the article is transcribed and no reference is redistributed. Reference adaptation: none was needed and none was made; every reference inside the delivered unit resolves inside it, so no unresolved reference or citation is produced. Printed slips kept literal: no printed word, symbol, hypothesis or number of the counted statement or of its proof was altered. Layout and class: the article's own class is \texttt{article} with packages including fullpage, amsfonts, amsmath, amssymb, amsthm, graphicx, sectsty, authblk; the wrapper is the lane's pinned \texttt{amsart} editorial wrapper at 10pt on A4, which restates the theorem-like environments the retained text needs and adds the article's own macro definitions that the retained text needs (none), with extra wrapper packages (bm, bbm, dsfont, xspace, cancel, mathtools). The delivered numbering is the unit's own. Assets: no retained span contains a figure environment, a graphics-inclusion command, a PDF-page inclusion, a table or a TikZ picture; no image file is copied into this package, the package declares no asset dependency and no image file is redistributed with it. Licence: version arXiv:2508.11102v1 of this article is distributed under the Creative Commons Attribution 4.0 International licence, as recorded in the arXiv licence metadata for this exact version and shown on the abstract page https://arxiv.org/abs/2508.11102v1. A full rights notice is delivered at licenses/arxiv-2508.11102-CC-BY-4.0.txt. Characters: every retained excerpt is pure ASCII, so the delivered engine, XeLaTeX with its default Latin Modern text font, reports no missing character.
\begin{eqnarray}\label{1.1}
-\frac{d^{2}}{dx^{2}}\psi_{j}(x)+q_{j}(x)\psi_{j}(0)=\lambda\psi_{j}(x),\,0<x<l_{j},\,j=1,2,\ldots,m,
\end{eqnarray}

\begin{equation}\label{1.2}
\psi_{j}(l_{j})=0,\,j=1,2,\ldots,m;\,\psi_{1}(0)=\psi_{2}(0)=\cdots=\psi_{m}(0),
\end{equation}

\begin{equation}\label{1.3}
\sum_{j=1}^{m}\big\{\frac{d}{dx}\psi_{j}(0)-\int_{0}^{l_{j}}\psi_{j}(x)\overline{q_{j}(x)}dx\big\}=0,
\end{equation}

\begin{eqnarray}\nonumber
\Phi(z):=\sum_{j=1}^{m}\Big(\int_{0}^{l_{j}}\sin{z(l_{j}-x)}\overline{q_{j}(x)}dx+\sin{z l_{j}}\sum_{n=1}^{\infty}\frac{q_{j,n}\overline{q_{j,n}}}{z^{2}-(\frac{n\pi}{l_{j}})^{2}}\Big)\prod_{k\neq j}\sin{z l_{k}}\\
+\sum_{j=1}^{m}\Big(z\cos{z l_{j}}+\sum_{n=1}^{\infty}\frac{(-1)^{n}n\pi}{l_{j}}\sin{z l_{j}}\frac{q_{j,n}}{z^{2}-(\frac{n\pi}{l_{j}})^{2}}\Big)\prod_{k\neq j}\sin{z l_{k}}=0,\label{Phi}
\end{eqnarray}

\begin{theorem}[Weyl's Law]\label{T31}
The zero density of $\Phi(z)$ is $\frac{\sum_{j=1}^{m}l_{j}}{\pi}$.
\end{theorem}

\begin{eqnarray}\nonumber
&&\delta(\sum_{j=1}^{m}\Big\{\sin{z l_{j}}\sum_{n=1}^{\infty}\frac{\widetilde{q}_{j,n}}{z^{2}-(\frac{n\pi}{l_{j}})^{2}}\Big\}\prod_{k\neq j}\sin{z l_{k}})\\\nonumber
&=&\max_{j}\{\delta(\Big\{\int_{0}^{l_{j}}\sin{z(l_{j}-x)}\widetilde{q}_{j}(x)dx\Big\}\prod_{k\neq j}\sin{z l_{k}})\}\\&=&\max_{j}\{\delta(\int_{0}^{l_{j}}\sin{z(l_{j}-x)}\widetilde{q}_{j}(x)dx)+\delta(\prod_{k\neq j}\sin{z l_{k}})\}\\
&=&\delta(\int_{0}^{l_{j_{0}}}\sin{z(l_{j_{0}}-x)}\widetilde{q}_{j_{0}}(x)dx)+\frac{\sum_{k\neq j_{0}}^{m}l_{k}}{\pi}\label{3.4},
\end{eqnarray}

\begin{eqnarray}\nonumber
&&\delta(\sum_{j=1}^{m}\int_{0}^{l_{j}}\sin{z(l_{j}-x)}\overline{\widetilde{q}_{j}(x)}dx\prod_{k\neq j}\sin{z l_{k}})\\
&=&\max_{j}\{\delta(\int_{0}^{l_{j}}\sin{z(l_{j}-x)}\overline{\widetilde{q}_{j}(x)}dx\prod_{k\neq j}\sin{z l_{k}})\}\\
&=&\max_{j_{0}}\{\delta(\int_{0}^{l_{j_{0}}}\sin{z(l_{j_{0}}-x)}\overline{\widetilde{q}_{j_{0}}(x)}dx)+\frac{\sum_{k\neq j_{0}}^{m}l_{k}}{\pi}\},\label{3.6}
\end{eqnarray}

\begin{eqnarray}\nonumber
&&\delta(\sum_{j=1}^{m}\sin{z l_{j}}\sum_{n=1}^{\infty}\frac{q^{1}_{j,n}\overline{q^{1}_{j,n}}-q^{2}_{j,n}\overline{q^{2}_{j,n}}}{z^{2}-(\frac{n\pi}{l_{j}})^{2}}\prod_{k\neq j}\sin{z l_{k}})\\
&=&\max_{j}\{\delta(\sin{z l_{j}}\sum_{n=1}^{\infty}\frac{q^{1}_{j,n}\overline{q^{1}_{j,n}}-q^{2}_{j,n}\overline{q^{2}_{j,n}}}{z^{2}-(\frac{n\pi}{l_{j}})^{2}})+\delta(\prod_{k\neq j}\sin{z l_{k}})\}\\
&=&\delta(\sin{z l_{j_{0}}}\sum_{n=1}^{\infty}\frac{q^{1}_{j_{0},n}\overline{q^{1}_{j_{0},n}}-q^{2}_{j_{0},n}\overline{q^{2}_{j_{0},n}}}{z^{2}-(\frac{n\pi}{l_{j_{0}}})^{2}})+\frac{\sum_{k\neq j_{0}}^{m}l_{k}}{\pi},\label{3.7}
\end{eqnarray}

\begin{eqnarray}\nonumber
&&\sum_{j=1}^{m}\Big(\int_{0}^{l_{j}}\sin{z(l_{j}-x)}\overline{\widetilde{q}_{j}(x)}dx+\sin{z l_{j}}\sum_{n=1}^{\infty}\frac{q^{1}_{j,n}\overline{q^{1}_{j,n}}-q^{2}_{j,n}\overline{q^{2}_{j,n}}}{z^{2}-(\frac{n\pi}{l_{j}})^{2}}\Big)\prod_{k\neq j}\sin{z l_{k}}\\
&\equiv&-\sum_{j=1}^{m}\Big\{\int_{0}^{l_{j}}\sin{z(l_{j}-x)}\widetilde{q}_{j}(x)dx\Big\}\prod_{k\neq j}\sin{z l_{k}},\,z\in\mathbb{C},\label{3.11}
\end{eqnarray}

\begin{theorem}[Partial Information]
Let the equations~(\ref{1.1}),~(\ref{1.2}), and~(\ref{1.3}) be defined on the fixed known star graph 
$T$ with $0$ as its vertex along with $m$ edges. We assume the convex hull of the support of $q_{j}(x)$ are strictly lesser than the length of edge 
$e_{j}$, $l_{j}$, $1\leq j\leq m$. If there are two sets of the potential functions $$\{q_{1}^{\sigma}(x),q_{2}^{\sigma}(x),\ldots,q_{m}^{\sigma}(x)\}\in L^{2}(0,l_{1})\oplus L^{2}(0,l_{2})\cdots\oplus L^{2}(0,l_{m}),\,\sigma=1,2,$$ that are imposed on the graph $T$ and sharing the identical $\Phi^{\sigma}(z),\,\sigma=1,2,$ constructed as in~(\ref{Phi}), then $q_{j}^{1}(x)\equiv q_{j}^{2}(x)$ almost everywhere for $1\leq j\leq m$.
\end{theorem}

\begin{proof}
Suppose we have two such set of potential functions  and assume that $\Phi^{1}(z)\equiv_{\mathbb{C}}\Phi^{2}(z)$.  This implies that they have the same zero density as given by Theorem \ref{T31}, which has the quantity
 $$\frac{\sum_{j=1}^{m}l_{j}}{\pi}.$$
Let us try to balance the zero densities on both sides of the equation~(\ref{3.11}), and  plug all of these common zeros into the following function
\begin{eqnarray}\nonumber
&&\sum_{j=1}^{m}\Big(\int_{0}^{l_{j}}\sin{z(l_{j}-x)}\overline{\widetilde{q}_{j}(x)}dx+\int_{0}^{l_{j}}\sin{z(l_{j}-x)}\widetilde{q}_{j}(x)dx+\int_{0}^{l_{j}}\sin{z(l_{j}-x)}G_{j}(x)dx\Big) \prod_{k\neq j}\sin{z l_{k}}.
\end{eqnarray}
Since for all $1\leq j\leq m$ by the theorem assumption,
\begin{eqnarray}
&&\delta(\int_{0}^{l_{j}}\sin{z(l_{j}-x)}\overline{\widetilde{q}_{j}(x)}dx)\,<\frac{l_{j}}{\pi};\\
&&\delta(\int_{0}^{l_{j}}\sin{z(l_{j}-x)}\widetilde{q}_{j}(x)dx)\,<\frac{l_{j}}{\pi};\\
&&\delta(\int_{0}^{l_{j}}\sin{z(l_{j}-x)}G_{j}(x)dx)<\frac{l_{j}}{\pi}.
\end{eqnarray}
Let us apply~(\ref{3.4}),~(\ref{3.6}), and~(\ref{3.7}) to reach the contradiction. The theorem is thus proven.

\end{proof}

\end{document}
